\documentclass[11pt]{article}
\usepackage{fontspec}
\usepackage[french]{babel}
\usepackage{amsmath}
\usepackage{amssymb}
\usepackage{amsthm}
\usepackage{mathrsfs}
\usepackage{tikz-cd}
\usepackage{enumitem}

\newcommand{\printedpage}[1]{\marginpar{\footnotesize p.~#1}}
\newcommand{\msp}[1]{\marginpar{\footnotesize [#1]}}
\newcommand{\Filt}{\mathrm{Filt}}
\newcommand{\Ob}{\mathrm{Ob}}
\newcommand{\Ss}{\mathrm{Ss}}
\newcommand{\M}{\mathcal{M}}
\newcommand{\Ci}{\mathcal{C}}

\begin{document}

\printedpage{34}
% head: A. Grothendieck : Dérivateurs, Ch. XV, Édition des Universités de Montpellier II et Paris VII
% foot: Version 5 novembre 2010
l'inclusion
\[
\underbrace{X'_{i_0}}_{\in\, \Filt^\pi} \;\longrightarrow\; X \;\simeq\; \varinjlim X_i \;.
\]
Comme $X'_{i_0} \in \Filt^\pi$, donc est $\pi$-accessible, il existe un $i_1 > i_0$ tel que $X'_{i_0} \longrightarrow X$ se factorise par
\[
X'_{i_0} \;\longrightarrow\; X_{i_1} \;.
\]
Appliquant ce même résultat à $X_{i_1}$, $X'_{i_1}$ \textit{etc.}, on trouve une suite infinie de morphismes
\[
\begin{tikzcd}[column sep=small]
X_{i_0} \arrow[dr] & & X_{i_1} \arrow[dr] & & X_{i_2} \arrow[dr] & \\
& X'_{i_0} \arrow[ur] & & X'_{i_1} \arrow[ur] & & X'_{i_2} \quad \cdots
\end{tikzcd}
\]
telle que les morphismes composés de deux en deux soient, sur la ligne supérieure, les transitions
\[
X_{i_m} \longrightarrow X_{i_n} \qquad (\; i_m \;<\; i_n \;) \;,
\]
et sur la ligne inférieure, les inclusions
\[
X'_{i_m} \hookrightarrow X'_{i_n} \qquad (\; \hookrightarrow X \;) \;.
\]
Soit $i_\omega = \sup_n i_n$ dans $I$, donc
\[
X_{i_\omega} \;=\; \varinjlim_{n} X_{i_n} \;\simeq\; \varinjlim_{n} X'_{i_n} \;\simeq\; X'_{i_\omega} \;,
\]
donc $X_{i_\omega} \longrightarrow X$ est un monomorphisme.

\msp{page 57}

La dernière assertion de 5.7 résulte aussitôt du fait qu'une limite inductive filtrante de monomorphismes dans $\M$ est un monomorphisme (condition $(e)$, p.~46).

On peut préciser la proposition 5.6 obtenue ainsi :

\textbf{Proposition 5.7}. \texttt{[Deux fois 5.7.]} \textit{Soient $\M$ une catégorie satisfaisant les conditions} (a), (b), (c), (d) (p.~44, 45) et (e) (p.~46). \textit{Soit $\M$ une partie de} $\Ob\M$\textit{, stable par isomorphie, $\pi$ un cardinal tel que les objets de} $\Filt^\pi$ \textit{soient $\pi$-accessibles. Les conditions suivantes sont équivalentes :} % [sic] Soit $\M$ (calligraphique) -> Soit $M$, la partie de Ob M est notée M partout ailleurs

\begin{enumerate}[label={}]
\item[a)] \textit{$M$ est $\pi$-accessible, et stable par $\varinjlim$ filtrantes.}
\item[b)] \textit{$M$ est $\pi$-accessible, et $M \cap \Filt^\pi$ est stable par $\varinjlim$ filtrante de cardinal $\leq \pi$.}
\item[c)] \textit{Il existe une partie $M^\pi$ de} $\Filt^\pi$\textit{, stable par $\varinjlim$ filtrantes de cardinal $\leq \pi$ et telle que pour tout $X \in \Ob\M$, on ait $X \in M$ si et seulement si l'ensemble $J_M$ des $Z_\alpha \in J$ tels que $Z_\alpha \in M^\pi$ est cofinal dans $J$.}
\end{enumerate}

(\textit{De plus, sous la condition} (c), \textit{on a}
\printedpage{35}
% head: A. Grothendieck : Dérivateurs, Ch. XV, Édition des Universités de Montpellier II et Paris VII
% foot: Version 5 novembre 2010
\msp{page 58}
$M^\pi = M \cap \Filt^\pi$.)
\begin{enumerate}[label={}]
\item[c$'$)] \textit{Il existe une partie $M_0$ de} $\Filt^\pi$ \textit{telle que $M$ soit la plus petite partie de $\M$ stable par limites inductives filtrantes.} (\textit{De plus, on peut alors prendre $M_0$} petite.)\footnote[22]{\textbf{N.B.} On a
\[
\begin{tikzcd}[ampersand replacement=\&]
(a) \arrow[r, Leftrightarrow] \arrow[d, Rightarrow] \& (b) \arrow[d, Rightarrow] \\
(c') \arrow[r, Rightarrow, "?"] \& (c)
\end{tikzcd}
\quad ,
\]
sans conditions sur $\pi$ (sauf $\pi \geq \pi_0$).}.
\end{enumerate}

\textbf{Corollaire 5.8.} \textit{Les applications}
\[
M \;\longmapsto\; M^\pi = M \cap \Filt^\pi \;,
\]
\[
M^\pi \quad\longmapsto\quad \parbox[t]{0.5\textwidth}{\textit{l'ensemble des $X \in \Ob\M$ tels que l'ensemble $\Ss_\pi^{M^\pi}(X)$ des $Z_\alpha \in \Ss_\pi(X)$ qui sont dans $M^\pi$ soit cofinal} (= \textit{plus petit sous-ensemble de $\Ob\M$ stable par $\varinjlim$ filtrantes})}
\]
\textit{établissent des bijections, inverses l'une de l'autre, entre l'ensemble des $M \subset \Ob\M$ satisfaisant les conditions équivalentes de la proposition} 5.7\textit{, et l'ensemble des parties $M^\pi$ de} $\Filt^\pi$ \textit{stables par $\varinjlim$ filtrantes de cardinal $\leq \pi$.}

\textsc{Démonstration de} 5.7. L'implication $(b) \Longrightarrow (c)$ est tautologique, compte tenu de la définition de la $\pi$-accessibilité de $M$, et de la proposition 5.3. Et $(a) \Longrightarrow (b)$ provient de la stabilité de $\Filt^\pi$ par $\varinjlim$ filtrantes de cardinal $\leq \pi$ (condition $(d)$).

Il reste donc à prouver $(c) \Longrightarrow (a)$. (C'est ici qu'on doit utiliser l'hypothèse particulière sur $\pi$.) Une fois ceci vu, on saura l'équivalence de $(a)$, $(b)$, $(c)$, et d'autre part, il est immédiat que $(a) \Longrightarrow (c')$, en
\msp{page 59}
prenant $M_0 = M \cap \Filt^\pi$. D'autre part, $(c') \Longrightarrow (c)$ (sans hypothèses sur $\pi\,$?), car on peut supposer que la partie $M_0$ de $\Filt^\pi$ est stable par $\varinjlim$ filtrantes de cardinal $\leq \pi$. Mais alors, connaissant l'équivalence de $(a)$, $(b)$ $(c)$, d'où aussitôt le corollaire, il s'ensuit que l'enveloppe stable par $\varinjlim$ filtrantes de $M_0$ dans $\M$ n'est autre que la partie accessible déduite de $M_0$ dont il est question dans $(c)$, et on gagne. % [sic] $(b)$ $(c)$ -> $(b)$, $(c)$, virgule manquante

Donc tout revient à prouver que, si tous les objets de $\Filt^\pi$ sont accessibles, alors $(c)$ implique $(a)$. Montrons d'abord que $(c)$ implique le critère d'appartenance (5.3) (p.~51), donc par la proposition 5.6 “revue et corrigée” dans la note 21, il s'ensuit déjà que $(c) \Longrightarrow (b)$, et il restera à établir $(b) \Longrightarrow (a)$, \textit{i.e.} que $(b)$ implique la stabilité de $M$ par $\varinjlim$ filtrantes quelconques (pas seulement grandes devant $\pi$).

Que $(c)$ implique (5.3) va résulter du lemme suivant, appliqué à $I = \Ss_\pi(X)$ :

\textbf{Lemme 5.8.} \texttt{[Deux fois 5.8.]} \textit{Soit $I$ un ensemble ordonné} $\pi$-adapté, i.e. \textit{filtrant et tel que toute partie filtrante $K$ de $I$, de cardinal $\leq \pi$, ait une borne supérieure $i_K$ dans $I$. Je dis qu'une partie $J$ de $I$ est} $\pi$-cofinale, \textit{si elle est cofinale, et si pour toute partie filtrante $K$ de $J$ de cardinal $\leq \pi$, sa borne supérieure $i_K$ dans $I$ est}
\printedpage{36}
% head: A. Grothendieck : Dérivateurs, Ch. XV, Édition des Universités de Montpellier II et Paris VII
% foot: Version 5 novembre 2010
\msp{page 60}
\textit{dans $J$.} (\textit{Il s'ensuit donc que $J$ est lui aussi $\pi$-adapté, pour la structure d'ordre induite.}) \textit{Ceci dit, l'intersection de deux parties $\pi$-cofinales de $I$ est encore $\pi$-cofinale.}

Il suffit de voir qu'elle est cofinale, car sa stabilité par bornes supérieures dans $I$ de parties filtrantes de cardinal $\leq \pi$ est tautologique (si elle est vraie pour des $J_\alpha$, elle est vraie pour $\bigcap_\alpha J_\alpha$). Soit $i_0 \in I$, et construisons un élément de $J \cap J'$ qui majore $i_0$. On construit de proche en proche une double suite
\[
\begin{tikzcd}
i_1 \arrow[dr] & & i_2 \arrow[dr] & & \cdots \\
& i'_1 \arrow[ur] & & i'_2 \arrow[ur] &
\end{tikzcd}
\]
avec les $i_\alpha$ dans $J$, les $i'_\alpha$ dans $J'$, et $i_1 \geq i_0$. On a donc
\[
\sup i_\alpha \;=\; \sup i'_\alpha \;,
\]
et si $i$ est la valeur commune des deux membres, on a $i \in J \cap J'$.

Ceci vu, pour prouver la validité de (5.3), on note que la partie $J' = \Ss_\pi^{(M^\pi)}(X)$ de $I = \Ss_\pi(X)$, formée des $Z_\alpha \subset X$ tels que $Z_\alpha \in M^\pi$, est stable par bornes supérieures de parties filtrantes de cardinal $\leq \pi$, vu que $M^\pi$ est stable
\msp{page 61}
par $\varinjlim$ de cardinal $\leq \pi$. Donc si $J$ est une partie $\pi$-cofinale de $\Ss_\pi(X)$, $J \cap J'$ est cofinale dans $J$, \textit{i.e.} dans $I$, si et seulement si $J'$ est cofinale dans $I$ (\textit{i.e.} si et seulement si $X \in M$, par la condition $(c)$). Car $J \cap J'$ cofinale $\Longrightarrow$ $J'$ cofinale tautologiquement, et $J'$ cofinale implique $J'$ $\pi$-cofinale, donc $J \cap J'$ $\pi$-cofinale par le lemme 5.8, donc cofinale.

Reste à prouver seulement que $(b) \Longrightarrow (a)$, \textit{i.e.} la stabilité de $M$ par limites inductives filtrantes quelconques, moyennant $(b)$. Ici à nouveau, cela ne doit pas utiliser l'hypothèse particulière sur $\pi$, ni même l'hypothèse d'exactitude $(e)$ sur $\M$, mais être une histoire générale d'accessibilité :

\textbf{Proposition 5.9.} \textit{Soit $\M$ une catégorie satisfaisant} (a), (b), (c), (d$_0$) (pages 44, 45), \textit{$\pi \geq \pi_0$ un cardinal, $M$ une partie $\pi$-accessible de} $\Ob\M$, i.e. \textit{satisfaisant la condition $b_\pi$ de} XIII, p.~88. \textit{Alors $M$ est stable par $\varinjlim$ filtrantes si et seulement si $M^\pi \overset{\text{déf}}{=} M \cap \Filt^\pi$ est stable par $\varinjlim$ filtrantes de cardinal $\leq \pi$.}

Je n'ai pas envie d'expliciter ici la démonstration, ayant l'impression pénible de refaire à perpette la même démonstration (ici et dans XIII). Je devrai revenir là-dessus dans une situation plus décantée.

\msp{page 62}

Je vais quand même esquisser la démonstration. Soit $X = \varinjlim\limits_{I} X_i$, $I$ filtrant, les $X_i$ dans $M$, il faut prouver que $X \in M$ (moyennant l'hypothèse de stabilité sur $M^\pi = M \cap \Filt^\pi$).
\printedpage{37}
% head: A. Grothendieck : Dérivateurs, Ch. XV, Édition des Universités de Montpellier II et Paris VII
% foot: Version 5 novembre 2010
\[
\begin{tikzcd}
& \Ci \arrow[r] \arrow[d] \arrow[dl, dotted, "\mathfrak{z}"'] \arrow[dr, phantom, "\text{\footnotesize cart.}"] & I \arrow[d, "\varphi"] \\
\underline{M}^\pi & \underline{M}^\pi/\underline{M} \arrow[l, dotted, "\text{\footnotesize “source”}"] \arrow[r, "\text{\footnotesize “but”}"'] & \underline{M}
\end{tikzcd}
\] % ?? uncertain: label of the dotted arrow -- printed glyph resembles a Fraktur z (\mathfrak{z}); possibly meant \exists

$\Ci$ est une catégorie cofibrée sur $I$, dont les fibres sont
\[
\Ci_i \;\simeq M^\pi/X_i \;,
\]
or comme $X_i \in M$, on voit, à cause de la condition de $\pi$-accessibilité, que $M^\pi/X_i$ est \textit{filtrante.} On va le dégager en

\textbf{Lemme 5.10.} \textit{Soit $M$ une partie $\pi$-accessible de $\M$. Alors pour tout $Y \in M$, $M^\pi/Y$ est une catégorie} filtrante, \textit{et $\Ss_\pi(Y)$ y est cofinale.}

\msp{page 63}

Considérons le foncteur canonique
\[
I \;=\; \Ss_\pi(Y) \xrightarrow{\;\varphi\;} \underline{M}^\pi/Y \;=\; I' \;,
\]
où la catégorie source $I$ est filtrante. Par SGA 4 I 8.1.3, ce foncteur est cofinal si et seulement s'il satisfait deux conditions F\,1, F\,2, et alors $I'$ est également filtrante. La condition F\,1 est la condition habituelle que tout $i' \in \Ob I'$ est “majoré” par un $\varphi(i)$, \textit{i.e.} tout $Z \longrightarrow Y$, avec $Z \in M^\pi$, est “majoré” par un \textit{monomorphisme} $Z_\alpha \hookrightarrow Y$, avec $Z_\alpha \in M^\pi$, ce qui résulte du fait que c'est majoré par $Z_\alpha \hookrightarrow Y$ avec $Z_\alpha \in \Filt^\pi$, \textit{i.e.} $Z_\alpha \in \Ss_\pi(Y)$, et que parmi ces $Z_\alpha$ ceux qui sont dans $M^\pi$ sont cofinaux (vu que $Y \in M$). L'autre condition\footnote[23]{Comme $I$ est filtrante.} est qu'une double flèche $i' \overset{f'}{\underset{g'}{\rightrightarrows}} \varphi(i)$ peut
\msp{page 64}
s'égaliser par une flèche $\varphi(h)$, avec $h : i \longrightarrow j$. Mais en fait, pour un $i' \in \Ob I'$, \textit{i.e.} un $Z \longrightarrow Y$ donné ($Z \in M^\pi$), et un $i \in \Ob I$, \textit{i.e.} un $Z_\alpha \hookrightarrow Y$ donné ($Z_\alpha \in M^\pi$), il existe au plus \textit{une} flèche dans $I' = \underline{M}^\pi/Y$ de $Z \longrightarrow Y$ dans $Z_\alpha \hookrightarrow Y$. Donc la condition F\,2 est trivialement satisfaite (on prend $h = \mathrm{id}_i$).

\textbf{Corollaire 5.11.} \textit{Si $(X_i)_{i \in I}$ est un système inductif filtrant, avec les $X_i$ dans $M^\pi$, alors la catégorie $\Ci = \underline{M}^\pi/\underline{M} \times_{\underline{M}} (I,\varphi)$, construite page} 62, \textit{est filtrante, et le foncteur canonique} (\textit{fibrant}) \textit{$\Ci \longrightarrow I$ est cofinal.}

Que $\Ci$ soit filtrante et $\Ci \longrightarrow I$ cofinal résulte formellement du fait que $I$ l'est, que les fibres de $\Ci$ le sont, et que $\Ci$ est cofibrée sur $I$. (Mérite un lemme, à toutes fins utiles.)

\textbf{N.B.} On peut remplacer $\Ci$ par $\Ci_0$ dont les fibres sont les
\[
\Ss_\pi^{M^\pi}(X_i) \;\subset\; \underline{M}^\pi \;; \qquad \text{cofinale dans } \Ci_i = \underline{M}^\pi/X_i \;,
\]
\msp{page 65}
on trouve une sous-catégorie de $\Ci$ cofinale dans $\Ci$, et qui est cette fois ordonnée tout comme $I$.

\printedpage{38}
% head: A. Grothendieck : Dérivateurs, Ch. XV, Édition des Universités de Montpellier II et Paris VII
% foot: Version 5 novembre 2010
On a
\[
X_i \;=\; \varinjlim_{Z_\alpha \in \Ci_i} Z_\alpha \;,
\]
donc
\[
X \;=\; \varinjlim_{\alpha \in \Ci} Z_\alpha \;,
\]
au total on est arrivé à représenter $X$ comme $\varinjlim$ filtrante d'éléments $Z_\alpha$ qui sont (non seulement dans $M$, mais même) dans $M^\pi$.

Changeons la notation, en prenant
\[
X \;=\; \varinjlim_{i \in I} X_i \;, \qquad I \text{ ensemble ordonné filtrant, } X_i \in M^\pi \;\; \forall\, i \in I \;,
\]
soit $I'$ l'ensemble des parties \textit{filtrantes} $J \subset I$ telles que $\mathrm{card}\, J \leq \pi$, $I'$ ordonné par inclusion. Alors $I'$ est grand devant $\pi$, et on a un système inductif
\[
(X_J)_{J \in I'} \;, \qquad X_J \;=\; \varinjlim_{i \in J} X_i \;.
\]
Par l'hypothèse de stabilité de $M^\pi$ (qui n'a pas encore servi), les $X_J$ sont dans $M^\pi$. On a donc
\[
X \;=\; \varinjlim_{J \in I'} X_J \;, \qquad \text{les } X_J \text{ dans } M^\pi\text{, } I' \text{ grande devant } \pi \;,
\]
donc $X \in M$ puisque $M$ est $\pi$-admissible, \quad \textit{q.e.d.}

\msp{page 66}

\textbf{Définition 5.12.}
\begin{enumerate}[label={}]
\item[1)] Soit $\M$ une catégorie satisfaisant les conditions $(a)$, $(b)$, $(c)$, $(d)$ (p.~44-46). Soit $\pi \geq \pi_0$ un cardinal. Une partie $M$ de $\Ob\M$ est dite \textit{L-$\pi$-accessible,} si elle est $\pi$-accessible et stable par petites $\varinjlim$ filtrantes\footnote[24]{On dit alors que $M$ est \textit{L-stable.}}, ou ce qui revient au même (5.9), si elle est $\pi$-accessible et si $M^\pi = M \cap \Filt^\pi$ est stable par $\varinjlim$ filtrantes de cardinal $\leq \pi$. (\textbf{N.B.} Si les objets de $\Filt^\pi$ sont $\pi$-accessibles, 5.8 nous dit que cela signifie aussi qu'il suffit pour cela qu'il existe une partie $M_0$ de $\Filt^\pi$ telle que $M$ soit la L-enveloppe (enveloppe stable par $\varinjlim$ filtrantes) de $M_0$ -- condition qui est de toutes façons toujours nécessaire, en prenant $M_0 = M^\pi$.)
\item[2)] Soit $M$ une partie de $\Ob\M$. On dit qu'elle est \textit{L-accessible,} si elle satisfait à l'une des quatre conditions équivalentes suivantes :
\begin{enumerate}[label={}]
\item[\textit{a})] Il existe un cardinal $\pi$ tel que $M$ soit L-$\pi$-accessible.
\item[\textit{a}$'$)] $M$ est accessible et L-stable.
\item[\textit{b})] Il existe une petite partie $M_0$ de $\Ob\M$, telle que $M$ soit la L-enveloppe de $M_0$\footnote[25]{Cette condition ne dépend pas de la filtration cardinale choisie de $\M$.}.
\item[\textit{b}$'$)] Il existe une partie $M_0$ de $\M$ contenue dans une partie $\Filt^\pi$ (ou ce qui revient au même, telle que la sous-catégorie pleine $\underline{M}_0$ de $\M$ soit essentiellement petite, \textit{i.e.} équivalente à une petite catégorie, telle que $M$ soit la \textit{L}-enveloppe de $M_0$).
\end{enumerate}
\end{enumerate}

\end{document}
